\section{Theoretical Results about Sparse Attention Mechanism}
\label{sec:theory}

In this section, we will show that that sparse attention mechanisms are as 
powerful and expressive as full-attention mechanisms in two respects.
First, we show that when sparse attention mechanisms are used in a standalone encoder 
(such as BERT), they are Universal Approximators of sequence 
to sequence functions in the style of \citet{Yun19}. 
We note that this property was also explored theoretically in contemporary work~\citet{yun2020on}.
%
Second, unlike~\citep{yun2020on}, we further show that sparse encoder-decoder transformers are Turing Complete 
(assuming the same conditions defined in~\citep{Perez19}).
%
Complementing the above positive results, we also show that moving to a sparse-attention mechanism 
incurs a cost, i.e.~there is no free lunch. In~\Cref{sec:limit}, we show lower bounds 
by exhibiting a natural task where any  sufficiently sparse  mechanism will 
require polynomially more layers. 

\subsection{Notation}
The complete Transformer {\em encoder} stack is nothing but the repeated application of a single-layer encoder (with independent parameters). 
We denote class of such Transformer encoders stack, defined using  generalized encoder (\Cref{sec:arch}),  by $\mathcal{T}_{D}^{H,m,q}$ which consists of $H$-heads with head size $m$ and $q$ is the hidden layer size of the output network, and the attention layer is defined by the directed graph $D$. 

The key difference between our proposed attention mechanism to that of \citet{vaswani2017attention,Yun19} is that we add a special token at the beginning of each sequence and assign it a special vector. 
We will refer to this as $\vx_0$. 
Therefore our graph $D$ will have vertex set $\set{0} \cup [n] = \set{0,1,2,\dots,n}$.
We will assume that this extra node and its respective vector will be dropped at the final output layer of transformer. 
To avoid cumbersome notation, we will still treat transformer as mapping sequences 
$\mX \in \R^{n \times d}$ to $\R^{n \times d}$. We will also allow the 
transformer to append position embeddings  $E \in \R^{d\times n}$ to 
matrix $X$ in the input layer. 

Finally, we need to define the function class and distance measure for proving universal approximation property. 
Let $\mathcal{F}_{CD}$ denote the set of continuous functions $f: [0,1]^{n \times d} \to \R^{n \times d} $ 
which are continuous with respect to the topology defined by  $\ell_p$ norm. 
Recall for any  $p \geq 1$, the $\ell_p$ distance is $d_p(f_1,f_2)  = \left(\int \norm{f_1(X) - f_2(X)}_p^p dX \right)^{1/p}$.

\subsection{Universal Approximators}
\begin{definition}
The star-graph $S$ centered at $0$ is the graph defined on $\set{0,\dots, n}$. 
The neighborhood of all vertices $i$ is $N(i) = \{0,i\} $  for $i \in \{1\dots n\}$ and 
$N(0) = \{1,\dots n\}$. 
\end{definition} 
Our main theorem is that the sparse attention mechanism defined by any graph containing $S$
is a universal approximator:
\begin{theorem}
\label{thm:universal}
    Given $1 < p < \infty$ and $\epsilon > 0 $, for any $f \in \mathcal{F}_{CD}$, there exists a 
    transformer with sparse-attention, $g \in \mathcal{T}_D^{H,m,q}$  such that 
    $d_p(f,g)\leq \epsilon$ where $D$ is any graph containing star graph $S$. 
\end{theorem}
To prove the theorem, we will follow the standard proof structure outlined in~\citep{Yun19}.

\textbf{Step 1: Approximate $\mathcal{F}_{CD}$ by piece-wise constant functions.} 
    Since $f$ is a continuous function with bounded domain $[0,1)^{n \times d}$, we will 
    approximate it with a  suitable piece-wise constant function. This is accomplished by a suitable partition of the  region $[0,1)$ into a grid of granularity $\delta$ to get a discrete set $\Gd$. Therefore, we can assume that we are dealing with a function 
    $\bar{f}: \Gd \to \R^{n \times d}$, where $d_p(f,\bar{f}) \leq \frac{\epsilon}{3}$. 
    
\textbf{Step 2:  Approximate piece-wise constant functions by modified transformers.} 
    This is the key step  of the proof where  the self-attention mechanism is used to generate a \emph{contextual-mapping} of the input. Informally, a contextual mapping is a unique code
    for the pair consisting of a matrix $(\mX,\vx_{i})$ and a column. Its uniqueness allows
    the Feed forward layers to use each code to map it to a unique output column. 
    
    The main technical challenge is computing the contextual mapping using only sparse attention mechanism. 
    This was done in \cite{Yun19} using a ``selective'' shift operator which shift up entries that are in a 
    specific interval. Key to their proof was the fact that the shift, was exactly the range of the largest 
    entry to the smallest entry.  
    
    Creating a contextual mapping with a sparse attention mechanism is quite a challenge.
    In particular, because each query only attends to a few keys, it is not at all clear
    that sufficient information can be corralled to make a contextual embedding of the 
    entire matrix. To get around this, we develop a sparse shift operator which shifts the entries of the matrices if they lie in a certain range. The exact amount of the shift is controlled by the directed sparse attention graphg $D$. The second key ingredient is the use of additional global token.  By carefully applying the operator
    to a set of chosen ranges, we will show that each column will contain a unique mapping of the
    full mapping.   Therefore, we can   
    augment the loss of inner-products in the self attention mechanism by using multiple 
    layers and an auxiliary global token. 
    
\textbf{Step 3: Approximate modified transformers by original Transformers}: The final 
step is to approximate the modified transformers by the original transformer which uses ReLU and 
softmax. 

We provide the full details in~\Cref{sec:apndx-universal}.

\subsection{Turing Completeness}
Transformers are a very general class. In the original paper of \citet{vaswani2017attention}, they were used 
in both an encoder and a decoder. While the previous section outlined how powerful just the encoders were,  another natural question is to ask what the
additional power of both a decoder along with an encoder is? \citet{Perez19} showed that the 
full  transformer based on a quadratic attention mechanism is Turing Complete. This result 
makes one  unrealistic assumption, which is that the model works on arbitrary precision 
model. Of course,  this is necessary as otherwise, Transformers are bounded finite 
state machines and cannot be Turing Complete. 

It is natural to ask if the full attention mechanism is necessary. Or can a 
sparse attention mechanism also be used to simulate any Turing Machine? 
We show that this is indeed the case: we can use a sparse encoder and sparse decoder
to simulate any Turing Machine. 

To use the sparse attention mechanism in the transformer architecture, we need to 
define a suitable modification where each token only reacts to previous tokens. 
Unlike the case for BERT, where the entire attention mechanism is applied once, in  full 
transformers, the sparse attention mechanism at decoder side is used token by token. 
Secondly the work of \citet{Perez19}, uses each token as a representation of the tape 
history and  uses the full attention to move and retrieve the correct tape symbol.
Most of the construction of \citet{Perez19} goes through for sparse attentions, except for their 
addressing scheme to point back in history (Lemma B.4 in \citep{Perez19}).
We show how to simulate this using a sparse attention mechanism and defer the 
details to~\Cref{sec:apndx-turing}.

\subsection{Limitations}
\label{sec:limit}
We demonstrate a natural task which can be solved by the full attention mechanism in $O(1)$-layers.
However, under standard complexity theoretic assumptions, this problem requires 
$\tilde{\Omega}(n)$-layers for any sparse attention layers with $\tilde{O}(n)$ edges (not just \bigb). (Here $\tilde{O}$ hides poly-logarthmic  factors). 
Consider the simple problem of finding the corresponding furthest vector 
for each vector in the given sequence of length $n$. Formally,

\textbf{Task 1.} \ Given $n$ unit vectors $\{u_1,\dots,u_n\}$, find $f(u_1,\dots,u_n) \to (u_{1^*}, \dots, u_{n^*})$  where for a fixed $j \in [n]$, we define $ j^* = \argmax_{k} \|u_k - u_j\|_2^2$.

Finding vectors that are furthest apart boils down to minimize 
inner product search in case of unit vectors. For a full-attention mechanism 
with appropriate query 
and keys, this task is very easy as we can evaluate all pair-wise inner products. 

The impossibility for sparse-attention follows from hardness results 
stemming from Orthogonal Vector Conjecture(OVC) \citep{abboud2014consequences,abboud2015tight,backurs2015edit,williams2005new}. The OVC  is a widely used assumption
in fine-grained complexity. Informally, it states that one cannot determine if the minimum inner product among 
$n$ boolean vectors is $0$ in subquadratic time. In~\Cref{sec:apndx-limit}, we show a 
reduction using  OVC to show that if a transformer $g \in \mathcal{T}_D^{H=1,m=2d,q=0}$ for 
any sparse directed graph $D$ can evaluate the Task $1$, it can solve the orthogonal vector problem.
\begin{proposition}
    There exists a single layer full self-attention $g\in\mathcal{T}^{H=1,m=2d,q=0}$ that can 
    evaluate Task 1,  i.e. $g(u_1,...,u_n) = [u_{1^*},\dots, u_{n^*}]$, but for any sparse-attention 
    graph $D$ with $\tilde{O}(n)$ edges (i.e.~inner product evaluations), would require $\tilde{\Omega}(n^{1-o(1)})$ layers.
\end{proposition}
\vspace{-3mm}
We give a formal proof of this fact in~\Cref{sec:apndx-limit}.